% !TEX root = main.tex
\section{The Exact Condition}
\label{sec:exact}

This section describes the deterministic hull and the coherence space of $T$ as two
polytopes of one graph, and gives an exact condition for a gap at the dual
of a tensor of two formulas.

\subsection{$\STAB$ and $\QSTAB$}
\label{sec:exact-stab}

For a finite graph $G$ with vertex set $V$, a \emph{clique} is a set of
pairwise adjacent vertices and a \emph{stable set} a set of pairwise
non-adjacent vertices. The \emph{stable set polytope} $\STAB(G)$ is $\Mix$
of the indicators of the stable sets, and the \emph{clique-constrained
polytope} $\QSTAB(G)$ is the set of $x : V \to \mathbb{R}_{\ge 0}$ with
$\sum_{v \in C} x_v \le 1$ for every clique $C$~\cite{gls1988}. A stable
set meets every clique in at most one vertex, so
$\STAB(G) \subseteq \QSTAB(G)$. We prove that the two polytopes describe
the coherence space and the deterministic hull of $T$.

\begin{theorem}[graph polytopes\leanref{td-eq-qstab}\leanref{recipe-eq-stab}]
\label{thm:qstab}
Let $G$ be the exclusivity graph of $T$. Then $P(T) = \QSTAB(G)$, and the
deterministic hull of $T$ is $\STAB(G)$.
\end{theorem}

\noindent
Every element of $P(X \multimap Y)$ is a mixture of certain states
(Proposition~\ref{prop:first-order}). Every certain state lies below
the indicator of the graph of a total function. So $P(T)$ is cut out by
the products of two such indicators. Their supports are the maximal
cliques of $G$, so $P(T) = \QSTAB(G)$. An element of $\QSTAB(G)$ with values in $\{0,1\}$
is the indicator of a stable set, so the deterministic hull is $\STAB(G)$.

So $T$ has a gap iff $\STAB(G) \ne \QSTAB(G)$. A graph is \emph{perfect}
if every induced subgraph has chromatic number equal to its largest
clique size. Then $\STAB(G) = \QSTAB(G)$ iff $G$ is
perfect~\cite{lovasz1972,chvatal1975}. The cycle of length $5$ of
Theorem~\ref{thm:pentagon} is an induced subgraph of $G$ that is not
perfect.

\subsection{The Clique Property and Perfection}
\label{sec:exact-perfect}

A formula $X$ has the \emph{clique property} if $P(X)$ is the set of
mixtures of the indicators of the cliques of its coherence graph $G_X$.
Data types, their negations, and the formulas $\mathbf{n} \multimap
\mathbf{m}$ (Proposition~\ref{prop:first-order}) have it. So does every
formula built from data types with $\mathbin{\&}$ and $\oplus$, since
$P(A \mathbin{\&} B)$ is the product of $P(A)$ and $P(B)$ and
$P(A \oplus B)$ is the convex hull of the copies of $P(A)$ and $P(B)$ on
the two parts of the web~\cite{danos-ehrhard2011}. Lean checks this for the
formulas built from the data type $\mathbf{1}$ with $\mathbin{\&}$ and
$\oplus$\leanref{recipe-code}. Their coherence graphs are the
\emph{cographs}, the graphs with no induced path on four
vertices~\cite{corneil1981}, with $\mathbin{\&}$ as join, which adds all
edges between the two parts, and $\oplus$ as disjoint union.

An \emph{induced $4$-cycle} of a formula $X$ is a quadruple of points
$a, b, c, d$ with $a \coh b \coh c \coh d \coh a$ and with $a, c$ and
$b, d$ incoherent. The \emph{strong product} $G \boxtimes H$ of two
graphs has the pairs of vertices as vertices, and two distinct pairs are
adjacent iff each coordinate is equal or adjacent. By
Definition~\ref{def:mall}(iv), the coherence graph of $X \otimes Y$ is
$G_X \boxtimes G_Y$. We prove the following exact condition.

\begin{theorem}[second order\leanref{closure}\leanref{recipe-core}\leanref{second-iff}]
\label{thm:second-exact}
Let $X$ and $Y$ have the clique property.
(i) $X \otimes Y$ has no gap. Moreover, $(X \otimes Y)^{\perp}$ has no gap
iff
$G_X \boxtimes G_Y$ is perfect.
(ii) If $G_X$ and $G_Y$ are cographs, then $(X \otimes Y)^{\perp}$ has no
gap iff $G_X$ or $G_Y$ has no induced $4$-cycle.
\end{theorem}

\noindent
If neither $X$ nor $Y$ has a gap, then neither do $X \otimes Y$,
$X \mathbin{\&} Y$, and $X \oplus Y$. Moreover, $P((X \otimes Y)^{\perp})$ is
cut out by the products of the certain states of $X$ and of $Y$. The Lean proof shows a bipolar theorem for downward-closed sets of
$\{0,1\}$-valued functions by hyperplane separation, as for the
completeness theorem of Section~\ref{sec:calculi}. With the
clique property these products are the indicators of the products
$K \times L$ of cliques, which are the maximal cliques of
$F = G_X \boxtimes G_Y$. So $P((X \otimes Y)^{\perp}) = \QSTAB(F)$, and its certain states are the
indicators of the stable sets of $F$. The deterministic hull is therefore $\STAB(F)$, which equals $\QSTAB(F)$ iff $F$ is
perfect. In (ii), the gap when both graphs
have an induced $4$-cycle follows from Theorem~\ref{thm:c4}, which holds for all
formulas. The closure statements, the description of
$P((X \otimes Y)^{\perp})$ as the orthogonal of the products of certain
states, and (ii) for $X$ and $Y$ built from $\mathbf{1}$ with
$\mathbin{\&}$ and $\oplus$ are formalized. The description as
$\QSTAB(F)$ is formalized at $T$ (Theorem~\ref{thm:qstab}), and the
equivalence with perfection in (i) is cited.

For $X = \mathbf{n} \multimap \mathbf{m}$ the coherence graph is complete
multipartite, with the inputs as the $n$ classes of $m$ points each and
two points adjacent iff they lie in different classes. It has an induced
$4$-cycle iff $n, m \ge 2$: two points in each of two classes form one,
for $n = 1$ the graph has no edges, and for $m = 1$ it is complete. So $T$ has a gap iff each of $X, Y, X', Y'$
has at least two elements, and the core has no gap\leanref{recipe-core}.
The gap when all four have at least two elements follows from Theorem~\ref{thm:pentagon}. The absence of a gap otherwise uses (ii) for
formulas not built from $\mathbf{1}$, through the transfer below, which is
not formalized.

\subsection{Four-Cycles: a Direct Proof}
\label{sec:exact-direct}

We prove the direction of (ii) without a gap by an explicit
decomposition, without perfect graph theory. If neither side of
$A \mathbin{\&} B$ is a clique, an incoherent pair in $A$ and one in $B$
form an induced $4$-cycle. So a formula built from $\mathbf{1}$ with
$\mathbin{\&}$ and $\oplus$ that has no induced $4$-cycle is isomorphic,
by a bijection of the webs that preserves coherence in both directions, to
one built from $\mathbf{1}$ with $\oplus$ and
$A \mapsto A \mathbin{\&} \mathbf{1}$\leanref{tree-of-noc4}. Its points
are the nodes of a rooted forest, and two points cohere iff one lies on the
path from a root to the other. So $G_X$ is the comparability graph of the
forest~\cite{wolk1962,golumbic1978}. For $v : |Y| \to \mathbb{R}_{\ge 0}$
let $\omega(v)$ be the largest total weight of a clique of $G_Y$. Then
$z$ lies in $P((X \otimes Y)^{\perp})$ iff
$\omega(\sum_{x \in \pi} z(x, \cdot)) \le 1$ for every path $\pi$ from a
root\leanref{mem-p-iff}. Every such $z$ is a mixture of
$\{0,1\}$-valued functions whose sums along the paths are indicators of
stable sets of $G_Y$\leanref{decompose}. The proof is an induction on the
forest, whose step at a root is an induction on $Y$. At $\oplus$ the two
components are decomposed separately. At $\mathbin{\&}$ a stable set lies
in one component, and is drawn there with probability given by $\omega$,
after a split of the weights below the root that keeps $\omega$ within its
level along every path\leanref{absorb}.

Lean checks (ii) in both directions for $X$ and $Y$ built from
$\mathbf{1}$ with $\mathbin{\&}$ and $\oplus$\leanref{second-iff}. For the
other formulas in (ii) with nonempty webs, the clique property makes a
coherence-preserving bijection onto such a formula preserve the coherence spaces. This transfer is not formalized. With the cited part of (i), the proof shows that the strong product of a
forest comparability graph and a cograph is perfect. This consequence is
not formalized. This extends the theorem that the strong product of two forest
comparability graphs, the graphs with no induced path on four vertices
and no induced $4$-cycle, is perfect~\cite{ravindra1978}. We did not find the extension in the literature.

\subsection{Beyond Second Order}
\label{sec:exact-beyond}

An occurrence of a formula $D$ in a formula $A$ is \emph{positive} if it
lies under an even number of negations and \emph{negative} otherwise, the
derived connectives $\parr$, $\oplus$, and $\multimap$ being unfolded. So
the left argument of $\multimap$ is a negative position. Beyond second
order, we find a gap even when no negatively occurring tensor has an
induced $4$-cycle in both factors.

\begin{proposition}[a gap without four-cycles\leanref{gap-f3}]
\label{prop:f3}
The formula $(\mathbf{2} \mathbin{\&} \mathbf{1}) \multimap
((\mathbf{2} \multimap \mathbf{2}) \otimes (\mathbf{2} \multimap
\mathbf{2}))$ has a gap. Every tensor occurring negatively in it has a
factor with at most three points, so none has an induced $4$-cycle in
both factors.
\end{proposition}

\noindent
The formula is $((\mathbf{2} \mathbin{\&} \mathbf{1}) \otimes T)^{\perp}$.
An element outside the deterministic hull has value $\tfrac12$ at five points of the
web. Its membership in the coherence space uses the bounds of Theorem~\ref{thm:mall} at three incoherent pairs and
at one point of $T$. The five points form an induced cycle of length $5$ in the
coherence graph, so by Theorem~\ref{thm:mall} every certain state is $1$
at no more than two of them. So every mixture has total weight at most $2$
there, against $\tfrac52$. The second statement of the proposition is
checked by inspection of the formula.
